\documentclass[10pt,a4paper]{amsart}
\usepackage[margin=19mm]{geometry}
\usepackage{amsmath,amssymb,mathtools}
\usepackage[scr=rsfs,cal=euler]{mathalfa}
\usepackage{xfrac}
\usepackage[unicode,hidelinks,bookmarks=false]{hyperref}
\newcommand{\ie}{i.e.\ }
\newcommand{\cf}{cf.\ }
\newcommand{\resp}{respectively\ }
\newcommand{\Subsection}{\subsection}
\providecommand{\nobreakdash}{\nobreak\hbox{-}}
\setlength{\emergencystretch}{2em}
\multlinegap0pt
\usepackage{bm}
\usepackage{bbm}
\usepackage{dsfont}
\usepackage{xspace}
\usepackage{cancel}
\usepackage{mathtools}
\usepackage{amsthm}
\makeatletter
\@ifundefined{theorem}{\newtheorem{theorem}{Theorem}}{}
\@ifundefined{proposition}{\newtheorem{proposition}[theorem]{Proposition}}{}
\makeatother
\title[Stuck Knots: Rigidity, Invariants, and Unsticking Distance]{Stuck Knots: Rigidity, Invariants, and Unsticking Distance}
\author{Ioannis Diamantis}
\date{Source: arXiv:2602.18129v1 (CC BY 4.0)}
\begin{document}
\maketitle

\noindent\emph{Source and licence.} Stuck Knots: Rigidity, Invariants, and Unsticking Distance. Version arXiv:2602.18129v1 (CC BY 4.0), distributed under the Creative Commons Attribution 4.0 International licence, as recorded in the arXiv licence metadata for this exact version and shown on the abstract page https://arxiv.org/abs/2602.18129v1. Full rights notice: licenses/arxiv-2602.18129-CC-BY-4.0.txt.\par\smallskip

\noindent\textit{Editorial scope.} Counted unit: the Proposition whose source label is \texttt{None} in the article (source0.tex lines 471--473, source label \texttt{None}), delivered together with the article's own complete original proof of it (source0.tex lines 475--495, one \texttt{proof} environment of 21 lines). No mathematics is written, repaired or re-derived: statement and proof are the article's own text line for line. Required context retained uncounted: none. Every definition, display and label of those lines is retained unchanged. Omitted from this package and not redistributed: 1--470, 474--474, 496--1202. Nothing inside a retained excerpt is omitted or altered: every retained body is delivered exactly as the article's lines are written, so no omission receipt of any kind accompanies this package. Citations: the retained excerpts carry no citation command at all, so no bibliography entry of the article is transcribed and no reference is redistributed. Reference adaptation: none was needed and none was made; every reference inside the delivered unit resolves inside it, so no unresolved reference or citation is produced. Printed slips kept literal: no printed word, symbol, hypothesis or number of the counted statement or of its proof was altered. Layout and class: the article's own class is \texttt{article} with packages including geometry, fontenc, inputenc, lmodern, microtype, graphicx, amsmath, amssymb, booktabs, enumitem, xcolor, float, hyperref, array; the wrapper is the lane's pinned \texttt{amsart} editorial wrapper at 10pt on A4, which restates the theorem-like environments the retained text needs and adds the article's own macro definitions that the retained text needs (none), with extra wrapper packages (bm, bbm, dsfont, xspace, cancel, mathtools). The delivered numbering is the unit's own. Assets: no retained span contains a figure environment, a graphics-inclusion command, a PDF-page inclusion, a table or a TikZ picture; no image file is copied into this package, the package declares no asset dependency and no image file is redistributed with it. Licence: version arXiv:2602.18129v1 of this article is distributed under the Creative Commons Attribution 4.0 International licence, as recorded in the arXiv licence metadata for this exact version and shown on the abstract page https://arxiv.org/abs/2602.18129v1. A full rights notice is delivered at licenses/arxiv-2602.18129-CC-BY-4.0.txt. Characters: every retained excerpt is pure ASCII, so the delivered engine, XeLaTeX with its default Latin Modern text font, reports no missing character.
\begin{proposition}
The relations above uniquely determine $P_R$ for all stuck knot diagrams.
\end{proposition}

\begin{proof}
The proof proceeds by induction on the pair $(N, s)$ under lexicographic ordering, where $N = c + s$ is the total number of crossings (the sum of classical crossings $c$ and rigid-height vertices $s$).

\begin{itemize}
    \item \textbf{Base Case:} If $N=0$, the diagram $D$ contains no crossings of any type. In this case, $D$ is an unlink. Its value is uniquely determined by the normalization $P_R(\bigcirc)=1$ and the disjoint union rule $P_R(L \sqcup \bigcirc) = \frac{a-a^{-1}}{z} P_R(L)$.

    \item \textbf{Inductive Step (Stuck crossings):} If $s > 0$, choose a rigid-height vertex $L_{\ast}^{\pm}$. Applying the rigid crossing relation expresses $P_R(L_{\ast}^{\pm})$ as a linear combination:
    \[
    P_R(L_{\ast}^{\pm}) = t\,P_R(L_0) + r^{\pm 1}P_R(L_{\pm}).
    \]
    The diagram $L_0$ has total crossings $N-1$, and the diagram $L_{\pm}$ has total crossings $N$ but with $s-1$ rigid-height vertices. In both cases, the resulting diagrams are strictly lower in the lexicographic order $(N, s)$.

    \item \textbf{Inductive Step (Classical crossings):} If $s = 0$ and $N > 0$, the diagram contains only classical crossings. We follow the standard unknotting algorithm: by choosing a base point and traversing the knot, we can express $P_R(D)$ via the classical skein relation 
    \[
    aP_R(L_+) - a^{-1}P_R(L_-) = zP_R(L_0)
    \]
    as a linear combination of an unknot and diagrams with strictly fewer crossings ($N-1$). Since $s=0$ for all diagrams in this step, they remain lower in the lexicographic order.
\end{itemize}

Since the lexicographic order on $\mathbb{N} \times \mathbb{N}$ is well-founded, the recursive process terminates in a finite number of steps. Thus, $P_R(D)$ is uniquely determined as a polynomial in $\mathbb{Z}[a^{\pm 1}, z^{\pm 1}, t, r^{\pm 1}]$.
\end{proof}

\end{document}
